\paragraph{极限 Stokes 泛函的测度实现与对应耗散}
在归一化逆塔积分与有限对应等变性的基础上，极限 Stokes 泛函不仅是圆柱顶次形式上的线性泛函，而且可由唯一的 Radon 概率测度实现；在对应算子可由有限支局部分支平均表示时，其函数层随之获得精确的 Markov 方差耗散公式。

\begin{theorem}[极限 Stokes 泛函的 Radon 测度实现]\label{thm:app-solenoid-stokes-radon-measure-realization}
沿用定理 \ref{thm:app-solenoid-fundamental-current-and-stokes} 与定理 \ref{thm:app-stokes-dirichlet-energy-multiscale-approx} 的设定。
令 \(g_n\) 为拉回度量链，\(\mathrm{vol}_n\) 为对应体积形式，并记
\[
\mathrm{Vol}_n:=\int_{M_n}\mathrm{vol}_n,
\qquad
\mu_n:=\frac{1}{\mathrm{Vol}_n}\,\mathrm{vol}_n.
\]
则存在唯一 Borel 概率测度 \(\mu_S\) 于
\(
S=\varprojlim(M_n,\pi_n)
\)
上，使得对每个 \(n\) 都有
\[
(p_n)_\#\mu_S=\mu_n.
\]
等价地，对任意连续圆柱函数 \(f\circ p_n\) 成立
\[
\int_S f\circ p_n\,\mathrm d\mu_S
=
\int_{M_n}f\,\mathrm d\mu_n.
\]

进一步，若 \(\omega=[f_n\,\mathrm{vol}_n]\in\Omega_{\mathrm{cyl}}^m(S)\) 为任意圆柱顶次形式，则
\[
\mathcal I_S(\omega)
=
\mathrm{Vol}_0\int_S f_n\circ p_n\,\mathrm d\mu_S,
\]
且右端与代表层 \(n\) 无关。
因此 \(\mathcal I_S\) 可唯一延拓到圆柱函数的 \(L^1(\mu_S)\) 完备化，并等同于相对于 \(\mathrm{Vol}_0\mu_S\) 的积分。
\end{theorem}
\begin{proof}
由于 \(g_{n+1}=\pi_n^\ast g_n\)，有 \(\mathrm{vol}_{n+1}=\pi_n^\ast\mathrm{vol}_n\)。
对 \(f\in C(M_n)\)，覆盖换元公式给出
\[
\int_{M_{n+1}} f\circ \pi_n\,\mathrm d\mu_{n+1}
=
\frac1{\mathrm{Vol}_{n+1}}\int_{M_{n+1}}\pi_n^\ast(f\,\mathrm{vol}_n)
=
\frac{d_n}{\mathrm{Vol}_{n+1}}\int_{M_n}f\,\mathrm{vol}_n.
\]
而 \(\mathrm{Vol}_{n+1}=d_n\mathrm{Vol}_n\)，故
\[
\int_{M_{n+1}} f\circ \pi_n\,\mathrm d\mu_{n+1}
=
\int_{M_n}f\,\mathrm d\mu_n.
\]
因此 \((\pi_n)_\#\mu_{n+1}=\mu_n\)。

记 \(\mathscr C_{\mathrm{cyl}}(S)\) 为 \(S\) 上连续圆柱函数所成的代数。
对 \(f\circ p_n\in\mathscr C_{\mathrm{cyl}}(S)\) 定义
\[
\Lambda(f\circ p_n):=\int_{M_n}f\,\mathrm d\mu_n.
\]
由上面的兼容性，\(\Lambda\) 良定义。
并且 \(\mathscr C_{\mathrm{cyl}}(S)\) 为包含常数、对共轭封闭并分离点的代数：若 \(x\neq y\in S\)，则存在某层 \(n\) 使 \(p_n(x)\neq p_n(y)\)，再由 \(M_n\) 的连续函数可分离两点，故有某个圆柱函数分离 \(x,y\)。
由 Stone--Weierstrass 定理，\(\mathscr C_{\mathrm{cyl}}(S)\) 在 \(C(S)\) 中稠密。
又 \(\Lambda\) 正且 \(\Lambda(1)=1\)，故它唯一连续延拓为 \(C(S)\) 上的正线性泛函。
由 Riesz--Markov 表示定理，存在唯一 Borel 概率测度 \(\mu_S\) 满足
\[
\Lambda(F)=\int_S F\,\mathrm d\mu_S
\qquad(F\in C(S)).
\]
对 \(F=f\circ p_n\) 代入即得 \((p_n)_\#\mu_S=\mu_n\)。

现在取 \(\omega=[f_n\,\mathrm{vol}_n]\in\Omega_{\mathrm{cyl}}^m(S)\)。
由圆柱相容性与 \(\mathrm{vol}_{n+1}=\pi_n^\ast\mathrm{vol}_n\)，可知 \(f_{n+1}=f_n\circ \pi_n\)，从而 \(f_n\circ p_n\) 与层 \(n\) 无关。
再由 \(\mathrm{Vol}_n=D_n\mathrm{Vol}_0\) 及 \(\mathcal I_S\) 的定义，
\[
\mathcal I_S(\omega)
=
\frac1{D_n}\int_{M_n}f_n\,\mathrm{vol}_n
=
\frac{\mathrm{Vol}_n}{D_n}\int_{M_n}f_n\,\mathrm d\mu_n
=
\mathrm{Vol}_0\int_S f_n\circ p_n\,\mathrm d\mu_S.
\]
故 \(\mathcal I_S\) 正是相对于有限测度 \(\mathrm{Vol}_0\mu_S\) 的积分。
其在 \(L^1(\mu_S)\) 圆柱完备化上的唯一延拓因此就是通常的 \(L^1\)-积分延拓。证毕。
\end{proof}

\begin{proposition}[对应算子的 Stokes 迹--Markov 结构与精确方差耗散]\label{prop:app-solenoid-stokes-markov-variance-dissipation}
沿用命题 \ref{prop:app-solenoid-stokes-equivariance-finite-correspondence} 与定理 \ref{thm:app-solenoid-stokes-radon-measure-realization} 的设定。
再假设对每个 \(n\) 都存在整数
\(
d_n:=\deg(s_n)=\deg(t_n)
\)
与一组局部可测入射分支 \(\psi_{n,1},\dots,\psi_{n,d_n}\)，使得对任意有界可测函数 \(f\) 有
\[
T_n(f\,\mathrm{vol}_n)=(P_n f)\,\mathrm{vol}_n,
\qquad
P_n f=\frac1{d_n}\sum_{i=1}^{d_n}f\circ\psi_{n,i},
\]
并且这些分支与逆塔结构相容，从而对任意有界可测 \(f\) 都满足
\[
P_{n+1}(f\circ\pi_n)=(P_n f)\circ\pi_n.
\]
则有：
\begin{enumerate}
  \item \(P_n\) 在 \(L^2(M_n,\mu_n)\) 上为 Markov 算子，且 \(\mu_n\) 为其平稳测度。
  \item 对任意 \(f\in L^2(M_n,\mu_n)\)，
  \[
  \|P_nf\|_{L^2(\mu_n)}\le \|f\|_{L^2(\mu_n)},
  \]
  并且成立精确耗散恒等式
  \[
  \|f\|_{L^2(\mu_n)}^2-\|P_nf\|_{L^2(\mu_n)}^2
  =
  \int_{M_n}\Bigl(P_n(f^2)-(P_nf)^2\Bigr)\,\mathrm d\mu_n
  \]
  \[
  =
  \frac1{2d_n^2}\int_{M_n}\sum_{i,j=1}^{d_n}
  \bigl(f\circ\psi_{n,i}-f\circ\psi_{n,j}\bigr)^2\,\mathrm d\mu_n.
  \]
  \item 在圆柱函数上定义
  \[
  P(f\circ p_n):=(P_n f)\circ p_n.
  \]
  则 \(P\) 良定义，并唯一延拓为 \(L^2(S,\mu_S)\) 上的 Markov 收缩算子；其平稳性、\(L^2\) 收缩性与上式的方差耗散恒等式在逆极限上仍然成立。
\end{enumerate}
\end{proposition}
\begin{proof}
由分支平均表达，\(P_n\) 显然保持非负性且 \(P_n\mathbf 1=\mathbf 1\)，故为 Markov 算子。
又由命题 \ref{prop:app-solenoid-stokes-equivariance-finite-correspondence} 与定理 \ref{thm:app-solenoid-stokes-radon-measure-realization}，
\[
\int_{M_n}P_n f\,\mathrm d\mu_n
=
\frac1{\mathrm{Vol}_n}\int_{M_n}T_n(f\,\mathrm{vol}_n)
=
\frac1{\mathrm{Vol}_n}\int_{M_n}f\,\mathrm{vol}_n
=
\int_{M_n}f\,\mathrm d\mu_n,
\]
故 \(\mu_n\) 为平稳测度。

对有界可测 \(f\)，Jensen 不等式给出点态估计
\[
(P_n f)^2\le P_n(f^2).
\]
两边对 \(\mu_n\) 积分并利用 \(f^2\) 的平稳性，即得
\[
\|P_n f\|_{L^2(\mu_n)}^2
=
\int_{M_n}(P_n f)^2\,\mathrm d\mu_n
\le
\int_{M_n}P_n(f^2)\,\mathrm d\mu_n
=
\int_{M_n}f^2\,\mathrm d\mu_n
=
\|f\|_{L^2(\mu_n)}^2.
\]
由有界函数在 \(L^2\) 中的稠密性，上式延拓到全部 \(L^2(M_n,\mu_n)\)。

再对固定点的有限平均应用恒等式
\[
\frac1{d_n}\sum_{i=1}^{d_n}a_i^2
-\left(\frac1{d_n}\sum_{i=1}^{d_n}a_i\right)^2
=
\frac1{2d_n^2}\sum_{i,j=1}^{d_n}(a_i-a_j)^2,
\]
取 \(a_i=f(\psi_{n,i}(x))\) 即得点态公式
\[
P_n(f^2)(x)-(P_nf(x))^2
=
\frac1{2d_n^2}\sum_{i,j=1}^{d_n}
\bigl(f(\psi_{n,i}(x))-f(\psi_{n,j}(x))\bigr)^2.
\]
由于右端关于分支对称，与局部编号无关，故可在全局上拼接。
对该式积分，并用 \(f^2\) 的平稳性，
\[
\|f\|_{L^2(\mu_n)}^2-\|P_nf\|_{L^2(\mu_n)}^2
=
\int_{M_n}\bigl(f^2-(P_nf)^2\bigr)\,\mathrm d\mu_n
=
\int_{M_n}\bigl(P_n(f^2)-(P_nf)^2\bigr)\,\mathrm d\mu_n,
\]
遂得到第二条。

对第三条，若圆柱函数 \(F\) 既可写成 \(f\circ p_n\) 又可写成 \(g\circ p_m\)，则在某共同提升层上二者一致。
由兼容性 \(P_{n+1}(f\circ\pi_n)=(P_n f)\circ\pi_n\)，定义
\(
P(f\circ p_n):=(P_n f)\circ p_n
\)
与代表无关，从而 \(P\) 在圆柱函数上良定义。
再由
\[
\int_S P(f\circ p_n)\,\mathrm d\mu_S
=
\int_{M_n}P_n f\,\mathrm d\mu_n
=
\int_{M_n}f\,\mathrm d\mu_n
=
\int_S f\circ p_n\,\mathrm d\mu_S
\]
以及层上的 \(L^2\) 收缩性，可知 \(P\) 在圆柱函数上平稳且收缩。
圆柱函数在 \(L^2(S,\mu_S)\) 中稠密，故 \(P\) 唯一延拓为 \(L^2(S,\mu_S)\) 上的 Markov 收缩算子。
耗散恒等式对圆柱函数由层上公式立即成立，再由 \(L^2\) 连续性延拓到逆极限上的全部 \(L^2\) 函数。证毕。
\end{proof}

\endinput
